\documentclass[11pt]{article}
\usepackage{amsmath}
\title{When a new stop authority changes a feasible retirement}
\date{26 November 2026}
\begin{document}
\maketitle
\paragraph{Declared authority and actuation model.}
The source replay and authored write-side gate establish why denying
a new write is not acknowledgement that an older one can have no
further effect. This study specifies what additional contract would
be needed to schedule prospective withdrawal of that older effect.
It does not obtain such a contract from the selected source code.

An older authorized action can start at any integer tick
$s\in[a,b]$. Its natural inclusive effect interval is $[s,s+h]$.
Its operator has a mandatory useful prefix through $s+m$ inclusive
with $0\le m\le h$. Effects in the prefix count as service only
as defined by this model; neither full mission completion nor
egress is inferred. A stop request sent at $x$ reaches the
actuator fence at $f\in[x+u,x+v]$. The fence blocks starts at
$f$ or later, including equality. For earlier starts, actuator
quiescence is guaranteed only from $f+g$, with $g\in[j,k]$.
All these intervals are closed, the time index is absolute,
and the uncertainties on $s,f,g$ combine freely. The most general
plant effect trace in the model includes every tick in
$[s,\min(s+h,f+g-1)]$ if $s<f$, and otherwise no effect.
A payload owner can revoke its permission from tick $C$ onward,
which imposes an obligation of no residual old effect there.
Any update to who can request the fence must itself be effective
before $x$. Authority is not inferred from delivery of a message,
and the fence is an additional enforceable actuator premise,
not a logical session retirement.

\paragraph{Exact timing.}
The guaranteed useful prefix requires both $s<f$ and
$s+m<f+g$. The worst cuts are latest start and earliest fence
arrival and quiescence, giving
$x\ge b+1-u$ and $x\ge b+m+1-u-j$.
Once those conditions hold, old effects can occur at or after
$C$ unless either they naturally end before $C$ or the
latest guaranteed transition to quiescence is no later than
$C$. These are precisely $b+h<C$ or $x\le C-v-k$.
For necessity when neither holds, take latest start, latest
fence arrival and latest settling lag; $\max(b,C)$ is a
realized offending effect tick. Thus the feasible send times
are an exact integer interval after also intersecting the
sender's effective authorization window. If relying on multiple
joint branches, compute these constraints within each branch
and intersect results; do not combine the earliest fence of
one branch with the latest physical effect of another.

Lean checks the quantified equivalence for one branch.
Independent explicit event histories and effect sets match
7{,}200 cases and 64{,}141 possible send choices. No report of
quiescence is generated merely by passing this model check.

\paragraph{Why the timing of new rights matters.}
Let the old action start at tick 1, naturally run through
tick 7, require its useful prefix through tick 4, and be forbidden
at tick 6 onward. Settling lag is one tick. A remote permitted
request arrives in 1--3 ticks; its required send would need to
be both at or after 3 and at or before 2. Granting a local
operator enforceable fence authority with arrival in 1--2 ticks
gives exactly one feasible send, tick 3. Making that right
effective at tick 4 is too late. If lag can instead be two ticks,
the same early local grant no longer suffices. A stock session
acknowledgement is not proof of the required lag bound.
Positive service in this example comes from changed lawful
powers with a better bounded transport contract, not from
crediting silence or broadening the owner's permission.
A direct planner with the same amended rights and mechanics
can reproduce the result.

\paragraph{Limit.}
Timer and authority readiness, tail bounds, execution observation
and freedom of branch combinations all need source justification.
The example does not show native lift quiescence, the permission
infrastructure, egress, a qualified controller, or a win over a
matched strong integrated stack. No contract disposition is made.
\end{document}
